\section*{Bab \ref{ch:g11:scal} --- Hasil Kali Skalar pada Bidang}

\begin{solution}{exo:g11:scal:1}
$(2,5)\cdot(3,-1) = 6 - 5 = 1$.

$(1,-4)\cdot(8,2) = 8 - 8 = 0$ (\omterm{def:g11:vect:vector}{vektornya} \omterm{def:g11:scal:orthogonal}{ortogonal}).

$\vec u\cdot\vec v = 3 \times 4 \times \cos\frac\pi3 = 12 \times \frac12
= 6$.
\end{solution}

\begin{solution}{exo:g11:scal:2}
$\vect{AB}\cdot\vect{AC} = 6 \times 6 \times \cos 60^\circ = 18$.

\omterm{def:g11:vect:vector}{Vektor} $\vect{AB}$ dan $\vect{BC}$ membentuk sudut
$180^\circ - 60^\circ = 120^\circ$ (letakkan pangkal keduanya di $B$), sehingga
$\vect{AB}\cdot\vect{BC} = 36 \cos 120^\circ = -18$.
\end{solution}

\begin{solution}{exo:g11:scal:3}
$\vec u \cdot \vec v = 3(t - 8) + t = 4t - 24 = 0$ memberi $t = 6$.

$\vec u - \vec v = (t - 4,\ 0)$, sehingga
$\vec u \cdot (\vec u - \vec v) = t(t-4) + 0 = t^2 - 4t$, yang lenyap
untuk $t = 0$ atau $t = 4$.
\end{solution}

\begin{solution}{exo:g11:scal:4}
$\vec u \cdot \vec v = 3 + 2 = 5$, $\norm{\vec u} = \sqrt5$,
$\norm{\vec v} = \sqrt{10}$, sehingga
\[
\cos\theta = \frac{5}{\sqrt5\,\sqrt{10}} = \frac{5}{5\sqrt2}
= \frac{\sqrt2}{2},
\]
jadi $\theta = 45^\circ$ tepat.
\end{solution}

\begin{solution}{exo:g11:scal:5}
$a^2 = 16 + 36 - 2 \times 4 \times 6 \times \cos\frac{2\pi}{3}
= 52 - 48 \times \left(-\frac12\right) = 76$, jadi $a = \sqrt{76} =
2\sqrt{19} \approx 8.7$.

Untuk segitiga kedua, aturan \omterm{def:g11:trigo:cossin}{kosinus} yang diselesaikan untuk sudutnya memberi
\[
\cos\widehat A = \frac{b^2 + c^2 - a^2}{2bc}
= \frac{25 + 16 - 49}{2 \times 5 \times 4} = -\frac{8}{40} = -\frac15,
\]
yang negatif: jadi $\widehat A$ tumpul.
\end{solution}

\begin{solution}{exo:g11:scal:6}
$\vect{AB}\,(4, -2)$ dan $\vect{AC}\,(2, 4)$:
$\vect{AB}\cdot\vect{AC} = 8 - 8 = 0$, sehingga sudut di $A$ siku-siku.
Sisinya: $AB = \sqrt{16 + 4} = \sqrt{20}$, $AC = \sqrt{4 + 16} =
\sqrt{20}$, dan $\vect{BC}\,(-2, 6)$ memberi $BC = \sqrt{40}$. Ketaatasasannya:
$BC^2 = 40 = 20 + 20 = AB^2 + AC^2$, yaitu aturan \omterm{def:g11:trigo:cossin}{kosinus} dengan
$\cos\widehat A = 0$ (Pythagoras).
\end{solution}

\begin{solution}{exo:g11:scal:7}
Bentuk pusat--jari-jarinya: $(x + 2)^2 + (y - 3)^2 = 25$.

Untuk lingkaran kedua, lengkapkan kuadratnya:
\[
x^2 - 6x + y^2 + 2y = 6
\iff (x - 3)^2 - 9 + (y + 1)^2 - 1 = 6
\iff (x-3)^2 + (y+1)^2 = 16 :
\]
yaitu berpusat $(3, -1)$ dan berjari-jari $4$.
\end{solution}

\begin{solution}{exo:g11:scal:8}
$M(x,y)$ berada pada lingkarannya ketika $\vect{MA}\cdot\vect{MB} = 0$ dengan
$\vect{MA}\,(-3 - x,\ -y)$ dan $\vect{MB}\,(-x,\ 2 - y)$:
\[
(-3 - x)(-x) + (-y)(2 - y) = x^2 + 3x + y^2 - 2y = 0 .
\]
Untuk titik asalnya: $0 + 0 + 0 - 0 = 0$, sehingga $O$ terletak pada lingkaran
itu --- secara geometris, $\vect{OA}\,(-3,0)$ dan $\vect{OB}\,(0,2)$ \omterm{def:g11:scal:orthogonal}{ortogonal},
sehingga $O$ memandang $[AB]$ dengan sudut siku-siku. (Setelah kuadratnya
dilengkapkan: pusatnya $\left(-\frac32, 1\right)$, jari-jarinya $\frac{\sqrt{13}}{2}$.)
\end{solution}

\begin{solution}{exo:g11:scal:9}
Garis tegak lurus melalui $P$ mempunyai \omterm{def:g11:vect:direction}{vektor arah} $d$, yaitu
$(-1, 2)$, sebagai \omterm{def:g11:vect:vector}{vektor} normalnya: \omterm{def:g10:algebra:equation}{persamaannya} $-x + 2y + c = 0$; melalui
$P(1,3)$: $-1 + 6 + c = 0$, jadi $c = -5$, sehingga $x - 2y + 5 = 0$.

Perpotongannya dengan $d$: dari $2x + y = 7$ diperoleh $y = 7 - 2x$; setelah
disubstitusikan, $x - 2(7 - 2x) + 5 = 5x - 9 = 0$, jadi $x = \frac95$,
$y = \frac{17}5$: yaitu $H\left(\frac95, \frac{17}5\right)$. Lalu
$\vect{PH}\left(\frac45, \frac25\right)$ dan
\[
PH = \sqrt{\frac{16}{25} + \frac{4}{25}} = \frac{\sqrt{20}}{5}
= \frac{2\sqrt5}{5} \approx 0.89 .
\]
\end{solution}

\begin{solution}{exo:g11:scal:10}
Menjumlahkan
\[
\norm{\vec u + \vec v}^2 = \norm{\vec u}^2 + 2\vec u\cdot\vec v +
\norm{\vec v}^2
\quad\text{dan}\quad
\norm{\vec u - \vec v}^2 = \norm{\vec u}^2 - 2\vec u\cdot\vec v +
\norm{\vec v}^2,
\]
suku silangnya saling menghapus dan
$\norm{\vec u + \vec v}^2 + \norm{\vec u - \vec v}^2 = 2\norm{\vec u}^2 +
2\norm{\vec v}^2$. Pada jajargenjang bersisi $\vec u$ dan $\vec v$,
diagonalnya adalah $\vec u + \vec v$ dan $\vec u - \vec v$: jadi jumlah
kuadrat kedua diagonalnya sama dengan jumlah kuadrat keempat
sisinya.
\end{solution}

\begin{solution}{exo:g11:scal:11}
\emph{1.} Karena $I$ \omterm{prop:g10:coordgeom:midpoint}{titik tengahnya}, $\vect{IB} = -\vect{IA}$ dan
$IA = 2$. Maka
\[
\vect{MA}\cdot\vect{MB}
= (\vect{MI} + \vect{IA})\cdot(\vect{MI} - \vect{IA})
= \norm{\vect{MI}}^2 - \norm{\vect{IA}}^2
= MI^2 - 4 .
\]

\emph{2.} Syarat $\vect{MA}\cdot\vect{MB} = 12$ menjadi
$MI^2 = 16$, yaitu $MI = 4$: jadi himpunannya adalah lingkaran berpusat $I$
dan berjari-jari $4$.
\end{solution}

\begin{solution}{pb:g11:scal:1}
\textbf{1.} $\vec u \cdot \vec v = -3 + 8 = 5$;
$\norm{\vec u} = 5$, $\norm{\vec v} = \sqrt5$;
$\cos\theta = \frac{5}{5\sqrt5} = \frac{1}{\sqrt5}$:
jadi $\theta \approx 63.4^\circ$.

\textbf{2.} $6 - 4k = 0$: $k = \frac32$.

\textbf{3.} $W = 50 \times 10 \times \cos 60^\circ = 250$
joule: hanya komponen sepanjang geraknya yang melakukan usaha.

\textbf{4.} $c^2 = 25 + 49 - 2 \times 35 \times \frac12 = 39$:
$c = \sqrt{39}$.

\textbf{5.} $\cos\gamma = \frac{16 + 36 - 64}{2 \times 4
\times 6} = -\frac14$: negatif, jadi sudutnya tumpul:
$\gamma \approx 104.5^\circ$.

\textbf{6.} Lewat \omterm{def:g10:coordgeom:system}{koordinat}: $\vec u \cdot \vec v = \cos a \cos b +
\sin a \sin b$. Lewat geometri: keduanya \omterm{def:g11:vect:vector}{vektor} satuan, sehingga
$\vec u \cdot \vec v = 1 \times 1 \times \cos(a - b)$.
Menyamakan kedua perhitungan bilangan yang sama itu:
$\cos(a - b) = \cos a \cos b + \sin a \sin b$.

\textbf{7.} $\cos(a + b) = \cos(a - (-b)) = \cos a \cos(-b) +
\sin a \sin(-b) = \cos a \cos b - \sin a \sin b$.

\textbf{8.} $\sin(a + b) = \cos\left(\frac\pi2 - a - b\right)
= \cos\left(\left(\frac\pi2 - a\right) - b\right)
= \cos\left(\frac\pi2 - a\right)\cos b +
\sin\left(\frac\pi2 - a\right)\sin b
= \sin a \cos b + \cos a \sin b$.

\textbf{9.} $\cos 15^\circ = \cos 45^\circ \cos 30^\circ +
\sin 45^\circ \sin 30^\circ = \frac{\sqrt2}{2} \cdot
\frac{\sqrt3}{2} + \frac{\sqrt2}{2} \cdot \frac12 =
\frac{\sqrt6 + \sqrt2}{4} \approx 0.966$. Lalu
$\sin 75^\circ = \cos 15^\circ$: nilai yang sama.

\textbf{10.} Dengan $b = a$ pada pertanyaan 7 dan 8:
$\cos 2a = \cos^2 a - \sin^2 a = 2\cos^2 a - 1 = 1 - 2\sin^2 a$
(dengan memakai $\cos^2 + \sin^2 = 1$), dan
$\sin 2a = 2 \sin a \cos a$.

\textbf{11.} $\sin^2 15^\circ = \frac{1 - \cos 30^\circ}{2}
= \frac{2 - \sqrt3}{4}$, sehingga
$\sin 15^\circ = \frac{\sqrt{2 - \sqrt3}}{2}$. Mengkuadratkan
$\frac{\sqrt6 - \sqrt2}{4}$:
$\frac{6 - 2\sqrt{12} + 2}{16} = \frac{8 - 4\sqrt3}{16} =
\frac{2 - \sqrt3}{4}$: bilangan yang sama dalam dua busana.

\textbf{12.} $\vect{AP}$ terurai menjadi bagian sepanjang $d$
(yang tak terlihat oleh $\vec n$) dan bagian yang tegak lurus $d$, yang
panjangnya adalah jarak yang dicari; mengalikannya secara titik dengan normal
satuan $\frac{\vec n}{\norm{\vec n}}$ mematikan bagian pertamanya dan
mengukur bagian keduanya. Dengan $A$ pada $d$:
$\vect{AP} \cdot \vec n = a x_0 + b y_0 - (a x_A + b y_A) =
a x_0 + b y_0 + c$, dan dari sanalah rumusnya. Contohnya:
$\frac{\abs{15 + 24 - 12}}{5} = \frac{27}{5} = 5.4$.

\textbf{13.} Jari-jarinya $= 5.4$:
$(x - 5)^2 + (y - 6)^2 = 29.16$.

\textbf{14.} Jaraknya $= \frac{\abs{-12}}{5} = 2.4$. Kakinya: garis
tegak lurus melalui titik asal adalah $(3t, 4t)$; pada garisnya:
$9t + 16t = 12$, jadi $t = \frac{12}{25}$: kakinya
$\left(\frac{36}{25}, \frac{48}{25}\right)$, pada jarak
$5t = \frac{60}{25} = 2.4$. Rumusnya dan perhitungan yang jujur
memberi hasil yang sama.

\textbf{15.} Alasnya $AB = 6$, tingginya $=$ jarak dari $C$ ke
$(AB)$ (yaitu sumbu $y = 0$): $5$; luasnya $15$. Garis $(AC)$:
$5x - 2y = 0$, dan
$\operatorname{dist}(B, (AC)) = \frac{\abs{30}}{\sqrt{29}}$;
lalu $\frac12 \times AC \times \frac{30}{\sqrt{29}} =
\frac12 \times \sqrt{29} \times \frac{30}{\sqrt{29}} = 15$:
luas yang sama, lewat garis tinggi yang mana pun.

\textbf{16.} $\operatorname{dist} = \frac{\abs{2 + 1 - 6}}
{\sqrt2} = \frac{3}{\sqrt2} \approx 2.12 > 2$: jadi garisnya lewat
tanpa menyentuh lingkarannya --- terpisah.

\textbf{17.} Sisipkan $I$:
$\vect{MA} \cdot \vect{MB} = (\vect{MI} + \vect{IA}) \cdot
(\vect{MI} + \vect{IB}) = MI^2 + \vect{MI} \cdot (\vect{IA} +
\vect{IB}) + \vect{IA} \cdot \vect{IB}$. Suku tengahnya mati
($\vect{IA} + \vect{IB} = \vec 0$), dan
$\vect{IA} \cdot \vect{IB} = -IA^2 = -\frac{AB^2}{4}$
(dua \omterm{def:g11:vect:vector}{vektor} berlawanan yang panjangnya $\frac{AB}2$).

\textbf{18.} $\vect{MA} \cdot \vect{MB} = 0 \iff MI^2 =
\frac{AB^2}{4} \iff MI = \frac{AB}{2}$: yaitu lingkaran berpusat
$I$ dan berjari-jari $\frac{AB}2$ --- lingkaran berdiameter
$[AB]$.

\textbf{19.} \omterm{def:g10:algebra:expand}{Menjabarkan} setiap pasangannya dari $H$ ($\vect{BC} =
\vect{HC} - \vect{HB}$, dan seterusnya):
\[
\vect{HA} \cdot \vect{HC} - \vect{HA} \cdot \vect{HB}
+ \vect{HB} \cdot \vect{HA} - \vect{HB} \cdot \vect{HC}
+ \vect{HC} \cdot \vect{HB} - \vect{HC} \cdot \vect{HA} = 0 :
\]
semuanya saling menghapus berpasangan --- jadi identitasnya berlaku untuk
\emph{sebarang} empat titik. Untuk $H$ kita, dua hasil kali titik yang
pertama lenyap menurut hipotesisnya, sehingga
$\vect{HC} \cdot \vect{AB} = 0$: jadi $H$ terletak pada garis tinggi dari
$C$. Tiga garis tinggi, satu titik.

\textbf{20.} Rumus penjumlahan sudut datang dari menghitung satu hasil kali
titik dengan wajah \emph{\omterm{def:g10:coordgeom:system}{koordinat}} dan wajah
\emph{norma-dan-sudut}; rumus jaraknya dari wajah
\emph{proyeksi}; \omterm{pb:g11:scal:1}{titik tingginya} dari wajah
\emph{kelinearan ganda} (jabarkan dari satu titik dengan Chasles). Hitunglah
besaran yang sama dua kali, samakan, dan sebuah teorema jatuh sendiri --- itulah
seluruh keahlian \omterm{def:g11:scal:dot}{hasil kali skalar}.

\end{solution}
